% Part: first-order-logic
% Chapter: natural-deduction
% Section: proving-things

\documentclass[../../../include/open-logic-section]{subfiles}

\begin{document}

\iftag{FOL}
      {\olfileid[tr]{fol}{ntd}{pro}}
      {\olfileid[tr]{pl}{ntd}{pro}}

\olsection{\usetoken{P}{derivation} Örnekleri}

\begin{ex}
!!a{derivation} verelim; bunun sonucu !!{sentence}
$(!A \land !B) \lif !A$ olacaktır.

Bu !!{derivation} için istenen sonucu alta yazarak başlarız.
\begin{prooftree}
\AxiomC{}
\UnaryInfC{$(!A\land !B) \lif !A$}
\end{prooftree}

Ardından, bu biçimde !!a{sentence} ile sonuçlanabilecek çıkarım türünü
belirlememiz gerekir. Sonuçta !!{main operator} olarak $\lif$ bulunduğundan
\Intro{\lif} kuralını kullanarak ulaşmayı deneriz. İlerlerken kullanılan
varsayımları ve çıkarım kurallarının etiketlerini yazmak en iyisidir; böylece
ispatın sonunda bütün varsayımların !!{discharged} olup olmadığı kolayca
görülebilir.
\begin{prooftree}
\AxiomC{$\Discharge{!A \land !B}{1}$}
\DeduceC{$!A$}
\DischargeRule{\Intro{\lif}}{1} 
\UnaryInfC{$(!A\land !B) \lif !A$}
\end{prooftree}

Şimdi $!A \land !B$ varsayımından $!A$'ya giden adımları doldurmamız
gerekir. Ele alınacak tek bağlaç $\land$ olduğundan $\land$ giderme kuralını
kullanmalıyız. Böylece şu ispatı elde ederiz:
\begin{prooftree}
\AxiomC{$\Discharge{!A \land !B}{1}$}
\RightLabel{\Elim{\land}}
\UnaryInfC{$!A$}
\DischargeRule{\Intro{\lif}}{1} 
\UnaryInfC{$(!A\land !B) \lif !A$}
\end{prooftree}
Artık $(!A \land !B) \lif !A$ için doğru bir !!{derivation} elde ettik.
\end{ex}

\begin{ex}
Şimdi $(\lnot !A \lor !B) \lif (!A \lif !B)$ için !!a{derivation}
verelim.

Bu !!{derivation} için istenen sonucu alta yazarak başlarız.
\begin{prooftree}
\AxiomC{}
\UnaryInfC{$(\lnot !A \lor !B) \lif (!A \lif !B)$}
\end{prooftree}
Bu sonucu verebilecek bir mantıksal kural bulmak için sonuçtaki mantıksal
bağlaçlara, yani $\lnot$, $\lor$ ve $\lif$ bağlaçlarına bakarız. Şimdilik
yalnızca $\lif$'in ilk geçişiyle ilgileniriz; çünkü son sekantta
!!{main operator} odur. Bu sekanttaki !!{sentence} içinde $\lnot$, $\lor$ ve
$\lif$'in ikinci geçişi başka bir bağlacın kapsamındadır; onlarla daha sonra
ilgileneceğiz.
Dolayısıyla \Intro{\lif} kuralıyla başlarız. Doğru bir uygulama şöyle
görünmelidir:
\begin{prooftree}
\AxiomC{$\Discharge{\lnot !A \lor !B}{1}$}
\DeduceC{$!A \lif !B$}
\DischargeRule{\Intro{\lif}}{1}
\UnaryInfC{$(\lnot !A \lor !B) \lif (!A \lif !B)$}
\end{prooftree}
Buradan devam etmenin iki yolu vardır. Ya aşağıdan yukarı çalışmayı sürdürüp
başka bir \Intro{\lif} uygulaması arayabiliriz ya da yukarıdan aşağı
çalışıp bir \Elim{\lor} kuralı uygulayabiliriz. İkinci yolu seçelim.
$\lnot !A \lor !B$ varsayımını \Elim{\lor} kuralının en soldaki öncülü
olarak kullanacağız. Geçerli bir \Elim{\lor} uygulamasında diğer iki öncül
sonuç $!A \lif !B$ ile özdeş olmalıdır; ancak bunların her biri sırasıyla
başka bir varsayımdan, yani $\lnot !A \lor !B$ ifadesindeki iki
bileşenden birinden türetilebilir. Böylece ortaya çıkan !!{derivation}
şöyle görünür:
\begin{prooftree}
\AxiomC{$\Discharge{\lnot !A \lor !B}{1}$}
\AxiomC{$\Discharge{\lnot !A}{2}$}
\DeduceC{$!A \lif !B$}
\AxiomC{$\Discharge{!B}{2}$}
\DeduceC{$!A \lif !B$}
\DischargeRule{\Elim{\lor}}{2}
\TrinaryInfC{$!A \lif !B$}
\DischargeRule{\Intro{\lif}}{1} 
\UnaryInfC{$(\lnot !A \lor !B) \lif (!A \lif !B)$}
\end{prooftree}

Sağdaki iki dalın her birinde, en uygun biçimde \Intro{\lif} kullanarak
$!A \lif !B$ ifadesini !!{derive} istiyoruz.
\begin{prooftree}
\AxiomC{$\Discharge{\lnot !A \lor !B}{1}$}
\AxiomC{$\Discharge{\lnot !A}{2}, \Discharge{!A}{3}$}
\DeduceC{$!B$}
\DischargeRule{\Intro{\lif}}{3}
\UnaryInfC{$!A \lif !B$}
\AxiomC{$\Discharge{!B}{2}, \Discharge{!A}{4}$}
\DeduceC{$!B$}
\DischargeRule{\Intro{\lif}}{4}
\UnaryInfC{$!A \lif !B$}
\DischargeRule{\Elim{\lor}}{2}
\TrinaryInfC{$!A \lif !B$}
\DischargeRule{\Intro{\lif}}{1} 
\UnaryInfC{$(\lnot !A \lor !B) \lif (!A \lif !B)$}
\end{prooftree}

Bu !!{derivation} içinde eksik kalan iki parça için, orta dalda $\lnot !A$
ile $!A$'dan ve diğer dalda $!A$ ile $!B$'den $!B$ için
!!{derivation}s gerekir. Önce ilkini ele alalım. $\lnot !A$ ve $!A$, \Elim{\lnot}
kuralının iki öncülüdür:
\begin{prooftree}
\AxiomC{$\Discharge{\lnot !A}{2}$}
\AxiomC{$\Discharge{!A}{3}$}
\RightLabel{\Elim{\lnot}}
\BinaryInfC{$\lfalse$}
\DeduceC{$!B$}
\end{prooftree}
\FalseInt kullanarak sonuç olarak $!B$'yi elde edip dalı tamamlayabiliriz.
\begin{prooftree}
\AxiomC{$\Discharge{\lnot !A \lor !B}{1}$}
\AxiomC{$\Discharge{\lnot !A}{2}$}
\AxiomC{$\Discharge{!A}{3}$}
\RightLabel{\Elim{\lnot}}
\BinaryInfC{$\lfalse$}
\RightLabel{\FalseInt}
\UnaryInfC{$!B$}
\DischargeRule{\Intro{\lif}}{3}
\UnaryInfC{$!A \lif !B$}
\AxiomC{$\Discharge{!B}{2}, \Discharge{!A}{4}$}
\DeduceC{$!B$}
\DischargeRule{\Intro{\lif}}{4}
\UnaryInfC{$!A \lif !B$}
\DischargeRule{\Elim{\lor}}{2}
\TrinaryInfC{$!A \lif !B$}
\DischargeRule{\Intro{\lif}}{1} 
\UnaryInfC{$(\lnot !A \lor !B) \lif (!A \lif !B)$}
\end{prooftree}

Şimdi en sağdaki dala bakalım. Burada !!{derivation} tanımının
\emph{varsayımların kapatılmasına izin verdiğini}, ama bunu
\emph{zorunlu kılmadığını} kavramak önemlidir. Başka bir deyişle, $!A$ ve
$!B$ varsayımlarından diğerini kullanmadan yalnızca birinden $!B$
türetebiliyorsak bunda sorun yoktur. $!B$'den $!B$'yi !!{derive} ise
çok kolaydır: $!B$ tek başına böyle !!a{derivation} olur ve hiçbir çıkarıma
gerek yoktur. Dolayısıyla $!A$ varsayımını silebiliriz.
\begin{prooftree}
\AxiomC{$\Discharge{\lnot !A \lor !B}{1}$}
\AxiomC{$\Discharge{\lnot !A}{2}$}
\AxiomC{$\Discharge{!A}{3}$}
\RightLabel{\Elim{\lnot}}
\BinaryInfC{$\lfalse$}
\RightLabel{\FalseInt}
\UnaryInfC{$!B$}
\DischargeRule{\Intro{\lif}}{3}
\UnaryInfC{$!A \lif !B$}
\AxiomC{$\Discharge{!B}{2}$}
\RightLabel{\Intro{\lif}}
\UnaryInfC{$!A \lif !B$}
\DischargeRule{\Elim{\lor}}{2}
\TrinaryInfC{$!A \lif !B$}
\DischargeRule{\Intro{\lif}}{1} 
\UnaryInfC{$(\lnot !A \lor !B) \lif (!A \lif !B)$}
\end{prooftree}
Tamamlanmış !!{derivation} içinde en sağdaki \Intro{\lif} çıkarımının
aslında hiçbir varsayımı kapatmadığına dikkat edin.
\end{ex}

\begin{ex}
Şimdiye dek \FalseCl{} kuralına ihtiyaç duymadık. Bu kural özeldir; çünkü
kural sonucunda bir alt-!!{formula} olmayan varsayımı kapatmamıza izin
verir. \FalseInt{} kuralıyla yakından ilişkilidir. Aslında \FalseInt{}
kuralı \FalseCl{} kuralının özel bir durumudur: yalnızca \FalseInt{}'a
izin veren ``sezgici mantık'' adlı bir mantık vardır. \FalseCl{} kuralı,
başka hiçbir şey işe yaramadığında başvurulacak son çaredir. Örneğin
$!A \lor \lnot !A$ ifadesini !!{derive} istediğimizi varsayalım.
Alışılmış stratejimiz, $\Intro{\lor}$ kullanarak $!A \lor \lnot !A$
ifadesini !!{derive} istemek olurdu. Ancak bunun için varsayım
olmaksızın ya $!A$'yı ya da $\lnot !A$'yı !!{derive} gerekirdi; bu
mümkün değildir. \FalseCl{} yardımımıza yetişir:
\begin{prooftree}
  \AxiomC{$\Discharge{\lnot(!A \lor \lnot !A)}{1}$}
  \DeduceC{$\lfalse$}
  \DischargeRule{\FalseCl}{1}
  \UnaryInfC{$!A \lor \lnot !A$}
\end{prooftree}
Şimdi $\lnot(!A \lor \lnot !A)$'dan $\lfalse$ için !!a{derivation}
arıyoruz. $\lfalse$, $\Elim{\lnot}$ kuralının sonucu olduğuna göre bunu
deneyebiliriz:
\begin{prooftree}
  \AxiomC{$\Discharge{\lnot(!A \lor \lnot !A)}{1}$}
  \DeduceC{$\lnot !A$}
  \AxiomC{$\Discharge{\lnot(!A \lor \lnot !A)}{1}$}
  \DeduceC{$!A$}
  \RightLabel{\Elim{\lnot}}
  \BinaryInfC{$\lfalse$}
  \DischargeRule{\FalseCl}{1}
  \UnaryInfC{$!A \lor \lnot !A$}
\end{prooftree}
$\lnot !A$ için !!a{derivation} bulma stratejimiz, $\Intro{\lnot}$
uygulamasını gerektirir:
\begin{prooftree}
  \AxiomC{$\Discharge{\lnot(!A \lor \lnot !A)}{1}, \Discharge{!A}{2}$}
  \DeduceC{$\lfalse$}
  \DischargeRule{\Intro{\lnot}}{2}
  \UnaryInfC{$\lnot !A$}
  \AxiomC{$\Discharge{\lnot(!A \lor \lnot !A)}{1}$}
  \DeduceC{$!A$}
  \RightLabel{\Elim{\lnot}}
  \BinaryInfC{$\lfalse$}
  \DischargeRule{\FalseCl}{1}
  \UnaryInfC{$!A \lor \lnot !A$}
\end{prooftree}
Burada $\lnot(!A \lor \lnot !A)$ varsayımına ve yeni varsayımımız $!A$'dan
$\Intro{\lor}$ yoluyla çıkan $!A \lor \lnot !A$ ifadesine $\Elim{\lnot}$
uygulayarak $\lfalse$'ı kolayca elde edebiliriz:
\begin{prooftree}
  \AxiomC{$\Discharge{\lnot(!A \lor \lnot !A)}{1}$}
  \AxiomC{$\Discharge{!A}{2}$}
  \RightLabel{\Intro{\lor}}
  \UnaryInfC{$!A \lor \lnot !A$}
  \RightLabel{\Elim{\lnot}}
  \BinaryInfC{$\lfalse$}
  \DischargeRule{\Intro{\lnot}}{2}
  \UnaryInfC{$\lnot !A$}
  \AxiomC{$\Discharge{\lnot(!A \lor \lnot !A)}{1}$}
  \DeduceC{$!A$}
  \RightLabel{\Elim{\lnot}}
  \BinaryInfC{$\lfalse$}
  \DischargeRule{\FalseCl}{1}
  \UnaryInfC{$!A \lor \lnot !A$}
\end{prooftree}
Sağ tarafta aynı stratejiyi kullanırız; tek fark, $!A$'yı \FalseCl ile elde
etmemizdir:
\begin{prooftree}
  \AxiomC{$\Discharge{\lnot(!A \lor \lnot !A)}{1}$}
  \AxiomC{$\Discharge{!A}{2}$}
  \RightLabel{\Intro{\lor}}
  \UnaryInfC{$!A \lor \lnot !A$}
  \RightLabel{\Elim{\lnot}}
  \BinaryInfC{$\lfalse$}
  \DischargeRule{\Intro{\lnot}}{2}
  \UnaryInfC{$\lnot !A$}
  \AxiomC{$\Discharge{\lnot(!A \lor \lnot !A)}{1}$}
  \AxiomC{$\Discharge{\lnot !A}{3}$}
  \RightLabel{\Intro{\lor}}
  \UnaryInfC{$!A \lor \lnot !A$}
  \RightLabel{\Elim{\lnot}}
  \BinaryInfC{$\lfalse$}
  \DischargeRule{\FalseCl}{3}
  \UnaryInfC{$!A$}
  \RightLabel{\Elim{\lnot}}
  \BinaryInfC{$\lfalse$}
  \DischargeRule{\FalseCl}{1}
  \UnaryInfC{$!A \lor \lnot !A$}
\end{prooftree}
\end{ex}

\begin{prob}
Aşağıdakileri gösteren !!{derivation}s verin:
\begin{enumerate}
\item $!A \land (!B \land !C) \Proves (!A \land !B) \land !C$.
\item $!A \lor (!B \lor !C) \Proves (!A \lor !B) \lor !C$.
\item $!A \lif (!B \lif !C) \Proves !B \lif (!A \lif !C)$.
\item $!A \Proves \lnot\lnot !A$.
\end{enumerate}
\end{prob}

\begin{prob}
Aşağıdakileri gösteren !!{derivation}s verin:
\begin{enumerate}
\item $(!A \lor !B) \lif !C \Proves !A \lif !C$.
\item $(!A \lif !C) \land (!B \lif !C) \Proves (!A \lor !B) \lif !C$.
\item $\Proves \lnot(!A \land \lnot !A)$.
\item $!B \lif !A \Proves \lnot !A \lif \lnot !B$.
\item $\Proves (!A \lif \lnot !A) \lif \lnot !A$.
\item $\Proves \lnot(!A \lif !B) \lif \lnot !B$.
\item $!A \lif !C \Proves \lnot (!A \land \lnot !C)$.
\item $!A \land \lnot !C \Proves \lnot (!A \lif !C)$.
\item $!A \lor !B, \lnot !B \Proves !A$.
\item $\lnot !A \lor \lnot !B \Proves \lnot(!A \land !B)$.
\item $\Proves (\lnot !A \land \lnot !B) \lif\lnot(!A \lor !B)$.
\item $\Proves \lnot(!A \lor !B) \lif (\lnot !A \land \lnot !B)$.
\end{enumerate}
\end{prob}

\begin{prob}
Aşağıdakileri gösteren !!{derivation}s verin:
\begin{enumerate}
\item $\lnot(!A \lif !B) \Proves !A$.
\item $\lnot(!A \land !B) \Proves \lnot !A \lor \lnot !B$.
\item $!A \lif !B \Proves \lnot !A \lor !B$.
\item $\Proves \lnot \lnot !A \lif !A$.
\item $!A \lif !B, \lnot !A \lif !B \Proves !B$.
\item $(!A \land !B) \lif !C \Proves (!A \lif !C) \lor (!B \lif !C)$.
\item $(!A \lif !B) \lif !A \Proves !A$.
\item $\Proves (!A \lif !B) \lor (!B \lif !C)$.
\end{enumerate}
(Bunların hepsi $\FalseCl$ kuralını gerektirir.)
\end{prob}

\end{document}
